\documentclass{beamer}
\usepackage[utf8]{inputenc}
\usepackage[dutch]{babel}
\usepackage{keystroke}
\usepackage{booktabs}
\usepackage{bbding} 
\usepackage{datetime}
\newdate{date}{20}{09}{2025}

\usepackage{hyperref}
% colored links
\hypersetup{
    colorlinks=true,
    linkcolor=blue,
    urlcolor=blue,
}

\usetheme{Madrid}
\usecolortheme{default}
% \setbeamercolor{frametitle}{bg=orange,fg=black}
\definecolor{NLLGGOrange}{RGB}{255,146,10}
\setbeamercolor{title}{bg=NLLGGOrange,fg=black}
\setbeamercolor{frametitle}{bg=NLLGGOrange,fg=black}
\setbeamercolor{structure}{fg=NLLGGOrange}

\newcommand{\light}[1]{\textcolor{gray}{#1}}
\newcommand{\reponame}{presentatie-fai}
\newcommand{\myfaiconfig}{https://github.com/slspeek/fai}
\newcommand{\myfaicmds}{https://github.com/slspeek/fai-cmds}
\newcommand{\mypreseed}{https://slspeek.github.io/debian/}

%------------------------------------------------------------
%This block of code defines the information to appear in the
%Title page
\title[Automatiseer installatie en configuratie] %optional
{Automatiseer je Linux installatie en configuratie met FAI}

\author[S.L.Speek] % (optional)
{Steven L. Speek}

\institute[NLLGG] % (optional)
{
  Nederlandse Linux Gebruikers Groep 
}

% \date[\today{}] {\today{}}
\date{\displaydate{date}}

\logo{\includegraphics[height=1cm]{img/nllgg-logo.jpg}}

%End of title page configuration block
%------------------------------------------------------------



%------------------------------------------------------------
%The next block of commands puts the table of contents at the 
%beginning of each section and highlights the current section:

% \AtBeginSection[]
% {
%   \begin{frame}<handout:0>
%     \frametitle{Inhoudsopgave}
%     \tableofcontents[currentsection]
%   \end{frame}
% }
\newif\iflattersubsect

\AtBeginSection[] {
    \begin{frame}<beamer>
      \frametitle{Inhoudsopgave}
      \tableofcontents[currentsection]
     \end{frame}
    \lattersubsectfalse
}

\AtBeginSubsection[] {
    \iflattersubsect
    \begin{frame}<beamer>
      \frametitle{Inhoudsopgave}
      \tableofcontents[currentsubsection]
    \end{frame}
    \fi
    \lattersubsecttrue
}
%------------------------------------------------------------


\begin{document}

%The next statement creates the title page.
\frame{\titlepage}


%---------------------------------------------------------
%This block of code is for the table of contents after
%the title page
\begin{frame}
\frametitle{Inhoudsopgave}
\tableofcontents
\end{frame}
%---------------------------------------------------------
\begin{frame}
\frametitle{Download de handout}
PDFs: \url{https://github.com/slspeek/\reponame/releases/latest}
\centering
\includegraphics[width=0.5\textwidth]{img/qr-release.png}
\end{frame}

\section{Over mij}


\begin{frame}
\frametitle{Steven L. Speek}
\begin{itemize}
  \item<2-> Actief bij \href{https://linux020.nl/}{Linux\textsuperscript{020}}
  \item<3-> Docent Linux bij \href{https://accscip.nl/}{Actief Computer Centrum Amsterdam}
  \item<4-> Debian gebruiker sinds Sarge (2005)
  \item<5-> Ontwikkelaar in Python en Java  
  \item<6-> Linux netwerkbeheerder
  \item<7-> Github: \url{https://github.com/slspeek}
  \item<8-> Hardloper: Volg mij op \href{https://www.strava.com/athletes/123214921}{Strava}
\end{itemize}
\end{frame}

\section{Probleemstelling}


\begin{frame}
\frametitle{Probleemstelling}
Na een Linux installatie ga je telkens een paar dingetjes aanpassen:
\begin{itemize}
  \item Upstream Firefox installeren
  \item Mattermost installeren
  \item Element installeren
  \item VLC installeren
  \item Gnome Tweaks installeren
  \item BASH prompt, met foutmeldingen, configureren
  \item Sneltoetsen instellen
  \item Mooie background instellen
  \item Enzovoorts
\end{itemize}
Dit kost tijd, en is foutgevoelig. Bovendien had je in die tijd uitleg kunnen geven over Linux.
\end{frame}

\section{Verschillende oplossingen}


\begin{frame}
\frametitle{Oplossingen}
Voor Debian gebaseerde systemen zijn er een paar oplossingen:
\begin{itemize}
  \item \href{https://wiki.debian.org/DebianInstaller/Preseed}{Preseed} 
  \begin{itemize}
    \item Eenvoudig
    \item Beperkt
    \item Alleen voor installatie
    \item \href{\mypreseed}{Mijn preseed bestanden}
  \end{itemize}
  \item \href{https://fai-project.org/}{FAI - Fully Automatic Installation}
  \begin{itemize}
     \item Zeer volwassen project (sinds 1999)
     \item Goed ontwerp
     \item \href{\myfaiconfig}{Mijn FAI configuratie}
     \item \href{\myfaicmds}{Mijn FAI commando's} 
  \end{itemize}
  \item \href{https://ansible.com/}{Ansible}
  \begin{itemize}
     \item Heel flexibel
     \item Niet specifiek voor Debian
     \item Niet specifiek voor installatie
     \item Vereist SSH toegang
  \end{itemize}
\end{itemize}
\end{frame}

\begin{frame}
\frametitle{Vergelijking van mijn Preseed en FAI}
\begin{table} % [htbp] suggests placement options: here, top, bottom, page
    \centering % Centers the table on the page
    \begin{tabular}{lccccc} % Defines column alignment: l=left, c=center, r=right
        \toprule % Top rule (from booktabs package)
                              & Mijn preseed & FAI \\
        \midrule % Mid rule (from booktabs package)
         \textbf{Automatisch} & \XSolid & \Checkmark \\
         \textbf{Multiboot} & \Checkmark & \XSolid \\
         \textbf{Gebruikersnaam} & \Checkmark & \XSolid \\
         \textbf{Live}  & \Checkmark & \Checkmark \\
         \textbf{Live persistentie} & \Checkmark & \XSolid \\
         \textbf{Geen Webserver nodig} & \XSolid & \Checkmark \\
          
        \bottomrule % Bottom rule (from booktabs package)
    \end{tabular}
\end{table}
\end{frame}

\begin{frame}
\frametitle{Motivatie voor FAI}
Nadelen van \href{\mypreseed}{Mijn Preseed}:
\begin{itemize}
  \item Webserver nodig
  \item Bestanden meegeven moeilijk
  \item Moeilijk uitbreidbaar
  \item Amateuristische herbouw van FAI
  \item Geen scheiding mechaniek en configuratie
  \item Geen i18n
  \item Live en installatie zijn niet hetzelfde
\end{itemize}
Voordelen van FAI:
\begin{itemize}
  \item Goed ontwerp
  \item Volwassen (sinds 1999)
  \item Actief project
  \item Veel voorbeelden
\end{itemize}
\end{frame}

\begin{frame}
\frametitle{Maanschot idee}
Draai FAI vanuit het late-command in een preseed.
\centering
\includegraphics[width=0.5\textwidth]{img/Cartoon_space_rocket.png}<1->
\end{frame}

\section{Fully Automatic Installation (FAI)}


\begin{frame}
\frametitle{Wat is FAI?}
\begin{itemize}
  \item Fully Automatic Installation
  \item Automatische installatie van Linux systemen
  \item \textbf{Debian}, Ubuntu, Mint, CentOS, Rocky Linux, AlmaLinux, ...
  \item Netwerk installatie
  \item \textbf{USB installatie}
  \item Chroot installatie
  \item \textbf{Live installatie}
  \item Bijwerken van met FAI gemaakte systemen
\end{itemize}
\end{frame}

\begin{frame}
\frametitle{FAI componenten}
\begin{itemize}
  \item FAI software
  \begin{itemize}
    \item Canoniek \url{https://fai-project.org}
    \item Persoonlijk \url{\myfaicmds} toegespitst op USB
     installaties en live systemen
  \end{itemize}
  \item FAI configuratie
  \begin{itemize}
    \item Canoniek \url{https://github.com/faiproject/fai-config}
    \item Persoonlijk \url{\myfaiconfig} met Nederlandse GNOME configuraties
  \end{itemize}
\end{itemize}
\end{frame}

\section{FAI Configuratie}


\begin{frame}
\frametitle{Zaken automatiseren in FAI}
Een \textbf{CLASS} kan bestaan uit:
\begin{itemize} 
  \item Verseisten dat bepaalde CLASSES zijn gedraaid
  \item Pakketten lijst
  \item Antwoorden op \href{https://wiki.debian.org/PackageManagement/Preseed}{debconf vragen}
  \item (Configuratie) bestanden
  \item Scripts
  \item<2-> Hooks
  \begin{itemize}
    \item Script dat voor een bepaalde FAI-fase wordt gedraaid
    \begin{itemize}
      \item Bijvoorbeeld \texttt{instsoft} om voordat de pakketten worden geïnstalleerd iets te doen
      \end{itemize}
  \end{itemize}
  \item<3-> \light{Lokale pakketten}
  \item<3-> \light{Variabelen}
  \item<3-> \light{Disklayout}
  \item<3-> \light{Tests}
\end{itemize}
\end{frame}

\begin{frame}
\frametitle{Configuratie nader bezien}
Configuratie is een map met:
\begin{itemize}
  \item \texttt{classes/} - Klasse variabelen en afhankelijkheden
  \item<2-> \texttt{package\_config/}  Pakketten lijsten
  \item<3-> \texttt{files/} - (Configuratie) bestanden
  \item<4-> \texttt{scripts/} - Scripts
  \item<5-> \texttt{hooks/} - Hooks
  \item<6-> \texttt{debconf/} - Debconf antwoorden
  \item<6-> \light{\texttt{disk\_config/} - Disklayouts}
  \item<6-> \light{\texttt{basefiles/} - Basis bestanden archief}
  \item<6-> \light{\texttt{pkgs/} - Lokale pakketten}
  \item<6-> \light{\texttt{tests/} - Tests}
\end{itemize}
\onslide<7->{Zie voor meer details de \href{https://fai-project.org/fai-guide/\#_a_id_c3_a_the_configuration_space_and_its_subdirectories}{FAI Guide over de configuratie}}
\end{frame}

\section{Voorbeelden van classes}


\begin{frame}[fragile]
\frametitle{Voorbeeld FANCY\_BASH\_PROMPT}
\textbf{Probleem:} 
Je wilt een mooie BASH prompt, die ook aangeeft of het laatste commando goed ging of niet
\begin{verbatim}
user@host:~$ (exit 42)
user@host:~(42)* true
user@host:~$
\end{verbatim}

 
\onslide<2->{\textbf{Conceptuele oplossing:} \texttt{/etc/bash.bashrc} moet worden aangepast}

\onslide<3->{\textbf{Oplossing:}}
\begin{itemize}
  \item<4-> Maak een bestand \href{\myfaiconfig/blob/main/files/etc/bash.bashrc/FANCY_BASH_PROMPT}{\texttt{files/etc/bash.bashrc/FANCY\_BASH\_PROMPT}} met de aanpassingen
  \item<5-> Maak een script \texttt{scripts/FANCY\_BASH\_PROMPT/10-copy-file} met:
  \begin{verbatim}#!/bin/bash
fcopy -M /etc/bash.bashrc
  \end{verbatim} 
   
\end{itemize}
\end{frame}

\begin{frame}[fragile]
\frametitle{Voorbeeld UNATTENDED\_UPGRADES}
 
\textbf{Probleem:} 
\begin{itemize}
  \item Pakket \texttt{unattended-upgrades} moet worden geïnstalleerd
  \item Onbeheerde updates moeten worden ingeschakeld
  \item Configuratie om van alle bronnen updates te accepteren
\end{itemize}

\onslide<2->{\textbf{Oplossing:}}
\begin{itemize}
  \item<2-> Maak een bestand \texttt{package\_config/UNATTENDED\_UPGRADES} met:
  \begin{verbatim}PACKAGES install
unattended-upgrades
  \end{verbatim} 
  \item<3-> Maak een bestand \texttt{debconf/UNATTENDED\_UPGRADES} met:
  \begin{verbatim}unattended-upgrades \
    unattended-upgrades/enable_auto_updates \
      boolean true
  \end{verbatim}
\end{itemize}
\end{frame}

\begin{frame}[fragile]
\frametitle{Voorbeeld UNATTENDED\_UPGRADES (vervolg)}
\onslide{\textbf{Oplossing (vervolg):}}
\begin{itemize}
  \item<2-> Maak een bestand \texttt{files/etc/apt/apt.conf.d/\symbol{92}} \texttt{52unattended-upgrades-local/UNATTENDED\_UPGRADES} met de configuratie
  \begin{verbatim}Unattended-Upgrade::Origins-Pattern {
        "origin=*";
 };
  \end{verbatim}
  \item<3-> Maak een script \texttt{scripts/UNATTENDED\_UPGRADES/10-configuration} met:
  \begin{verbatim}#!/bin/bash
fcopy /apt/apt.conf.d/52unattended-upgrades-local
  \end{verbatim}
\end{itemize}
\end{frame}

\begin{frame}[fragile]
\frametitle{Voorbeeld MATTERMOST}
\textbf{Probleem:} 
Je wilt Mattermost installeren, maar het zit niet standaard in Debian. Gelukkig heeft \texttt{extrepo} een bron hiervoor

\onslide<2->{\textbf{Conceptuele oplossing:}}
\begin{itemize}
  \item<3-> Installeer \texttt{extrepo}
  \item<4-> Voeg de Mattermost bron toe
  \item<5-> Installeer het pakket
\end{itemize}
\end{frame}
\begin{frame}[fragile]
\frametitle{Voorbeeld MATTERMOST (vervolg)}
\textbf{Concrete oplossing:}
\begin{itemize}
  \item<2-> Maak een bestand \texttt{class/MATTERMOST.deps} met:
  \begin{verbatim}EXTREPO
  \end{verbatim}
  \item<3-> Maak een bestand \texttt{hooks/instsoft.MATTERMOST} met de inhoud:
  \begin{verbatim}#! /bin/bash
$ROOTCMD extrepo enable mattermost
$ROOTCMD apt-get update
  \end{verbatim}
  \item<4-> Maak een bestand \texttt{package\_config/MATTERMOST} met:
  \begin{verbatim}PACKAGES install
mattermost-desktop
\end{verbatim} 
  \end{itemize}
\end{frame}

\begin{frame}
\frametitle{Volgorde van classes}
\begin{itemize}
  \item<2-> Eerst de afhankelijkheden   
  \item<3-> Bij variabelen geldt dat de waarde gezet in de laatste class wint
  \item<4-> Bij scripts en hooks geldt dat alle scripts worden uitgevoerd in de volgorde van de classes
  \item<5-> Bij meerdere files van verschillende classes, bepaalt de volgorde welke file het wordt. De laatste wint. Zie \href{https://fai-project.org/doc/man/fcopy.html}{\texttt{man fcopy}}
\end{itemize}
\end{frame}

\begin{frame}
\frametitle{Abstract voorbeeld}
Een FAI installatie met de classes: \texttt{CLASS\_A CLASS\_B CLASS\_C}, waarbij \texttt{CLASS\_B} 
afhankelijk is van \texttt{CLASS\_B\_1} en \texttt{CLASS\_B\_2}

\begin{itemize}
  \item<2-> \textsc{defclass}
  \begin{itemize}
    \item leest \texttt{classes/CLASS\_B.deps} en voegt \texttt{CLASS\_B\_1} en \texttt{CLASS\_B\_2} voor \texttt{CLASS\_B} in
    \item resultaat: \texttt{CLASS\_A CLASS\_B\_1 CLASS\_B\_2 CLASS\_B CLASS\_C}
  \end{itemize}
  \item<3-> \textsc{defvar}
  \begin{itemize}
    \item source \texttt{classes/CLASS\_A.var}
    \item source \texttt{classes/CLASS\_B\_1.var}
    \item source \texttt{classes/CLASS\_B\_2.var}
    \item source \texttt{classes/CLASS\_B.var}
    \item source \texttt{classes/CLASS\_C.var}
  \end{itemize}
  \item<4-> \textsc{instsoft}
  \begin{itemize}
    \item pakketten uit \texttt{packages\_config/} worden ge\"{i}nstalleerd voor
     \texttt{CLASS\_A CLASS\_B\_1 CLASS\_B\_2 CLASS\_B CLASS\_C}
  \end{itemize}
\end{itemize}
\end{frame}

\begin{frame}
\frametitle{Abstract voorbeeld (vervolg)}
\begin{itemize}
  \item \textsc{configure}
  \begin{itemize}
    \item scripts uit \texttt{scripts/CLASS\_A/} worden uitgevoerd
    \item scripts uit \texttt{scripts/CLASS\_B\_1/} worden uitgevoerd
    \item scripts uit \texttt{scripts/CLASS\_B\_2/} worden uitgevoerd
    \item scripts uit \texttt{scripts/CLASS\_B/} worden uitgevoerd
    \item scripts uit \texttt{scripts/CLASS\_C/} worden uitgevoerd
  \end{itemize}
  Bij \texttt{fcopy} commando's geldt dat de file van de laatste class wint.
\end{itemize}
\end{frame}


\section{Profielen}


\begin{frame}
\frametitle{Profielen}
Een profiel heeft:
\begin{itemize}
  \item<2-> Naam
  \item<3-> Beschrijving
  \item<4-> Een \textbf{geordende} lijst van \textbf{classes}
\end{itemize}

\onslide<5->{Profielen kunnen zijn voor:}
  \begin{itemize}
    \item<5-> Desktop
    \item<5-> Server
    \item<5-> Laptop
    \item<5-> Specifiek voor een gebruiker
    \item<5-> Specifiek voor een groep gebruikers
  \end{itemize}
\end{frame}

\begin{frame}[fragile]
\frametitle{Voorbeeld profiel `Gnome developer'}
Inhoud van het bestand \texttt{class/gnome-developer.profile}:
\begin{verbatim}
Name: Gnome developer
Description: Gnome desktop met developer tools
(...)
Classes: INSTALL FAIBASE DEBIAN DUTCH CONSOLE_TWEAKS \
  ADVANCED_CLI_TOOLS TUX GNOME_CORE BACKGROUND \
  FIREFOX DEVELOPER 
\end{verbatim}
\end{frame}

\section{Mijn FAI repositories}
\begin{frame}
\frametitle{Mijn FAI repositories}
\begin{itemize}
  \item \href{\myfaiconfig/}{Mijn FAI configuratie}
  \begin{itemize}
    \item Profielen voor:
    \begin{itemize}
      \item GNOME minimaal
      \item GNOME developer
      \item GNOME compleet
      \item XFCE minimaal
      \item Simpel
    \end{itemize}
    \item Nederlandse locale
    \item Eenvoudiger dan \url{https://github.com/faiproject/fai-config} (de canonieke FAI configuratie)
  \end{itemize}
  \item \href{\myfaicmds/}{Mijn FAI commando's}
  \begin{itemize}
    \item \texttt{Makefile} om snel aan de slag te kunnen
    \begin{itemize}
      \item Installeer FAI
      \item Maak een NFS root (het installatie live systeem)
      \item Maak een installatie ISO (eventueel met mirror)
      \item Maak een live ISO voor een profiel
      \item Start een VM met een ISO
    \end{itemize}
  \end{itemize}
\end{itemize}
Gebruik de instructies in de \href{\myfaicmds/blob/main/README.md}{README.md}
\end{frame}

\begin{frame}
\frametitle{Verder}
\begin{itemize}
  \item \href{https://fai-project.org/}{FAI project}
  \item \href{https://fai-project.org/fai-guide/}{FAI Guide}
  \item \href{\myfaiconfig/}{Mijn FAI configuratie} i.h.b. 
  de class \texttt{GNOME\_CORE} voor de \texttt{dconf} configuratie
  \begin{itemize}
    \item \href{\myfaiconfig/blob/main/files/etc/dconf/db/local.d/00-favorite-apps/DEVELOPER}{\texttt{dconf} configuratie voor de favorieten}
    \item \href{\myfaiconfig/blob/main/files/etc/dconf/db/local.d/10-shortcuts/GNOME_CORE}{\texttt{dconf} configuratie voor de sneltoetsen}
    \item \href{\myfaiconfig/blob/main/files/etc/dconf/db/local.d/40-gnome-extensions/GNOME_EXTENSIONS}{\texttt{dconf} configuratie voor de extensies}
  \end{itemize}
  \item Ga aan de slag met de instructies uit de \href{\myfaicmds/blob/main/README.md}{README.md}
\end{itemize}
\end{frame}

\section{Demo}
\begin{frame}
\frametitle{Demo}
\begin{itemize}
  \item Demo van FAI \texttt{live-simpel.iso} in \texttt{libvirt}(VM) (\href{https://youtu.be/4j1Wtoy2kK8}{Bekijk op YouTube})
  \item Installatie van \texttt{gnome-minimaal-demo} profiel (\href{https://youtu.be/CE_PjYTGtjc}{Bekijk op YouTube})
\end{itemize}
\end{frame}

\begin{frame}
\frametitle{Over deze presentatie}
\begin{itemize}
  \item Gemaakt met \href{https://nl.mirrors.cicku.me/ctan/macros/latex/contrib/beamer/doc/beameruserguide.pdf}{\LaTeX{} beamer}
  \item Broncode: \url{https://github.com/slspeek/\reponame}
  \item PDFs: \url{https://github.com/slspeek/\reponame/releases/latest}
  \item Experimenteren met \LaTeX{} beamer: \url{https://github.com/slspeek/empty-latex-beamer}
  \begin{itemize}
    \item NLLGG kleuren en logo
    \item Productie van PDFs via \LaTeX{} dockerimage
    \item Github Actions workflow met releases
  \end{itemize}
\end{itemize}
\end{frame}

\section{Vragen}
\begin{frame}
\frametitle{Vragen?}
\begin{itemize}
  \item Vragen?
  \item<2-> Opmerkingen?
  \item<3-> Verbeteringen?
\end{itemize}
\end{frame}

\begin{frame}
\frametitle{Volg mij op Strava}
\centering
\includegraphics[width=0.5\textwidth]{img/qr-follow-me-on-strava.png}
\end{frame}

\end{document}